\renewcommand{\theposttitle}{Fluid vs Crystalized Intelligence}
\renewcommand{\thepostdate}{2023-03-01}
\vspace*{\fill}
\begin{center}
\includegraphics[width=\textwidth, height=0.55\textheight, keepaspectratio]{C:/Users/James/Documents/GitHub/VisualIdeasBook/data/images/5d6dff56e22281550173223d586e5cbc.png}
\end{center}
\vspace{2mm}
\begin{center}
\begin{minipage}{0.92\textwidth}
\small
I've been thinking about age a lot recently so I was intrigued to come across the above concept from Catell's Intelligence Theory. Apparently, many physicists, scientists, and even entrepreneurs hit their innovation peak early in their careers. And they waste the latter end of their lives miserable because they continue to chase success to rival that of their youth. Turns out, there's a scientific reason for this inability.

There're two types of intelligence (and they're affected by age).

\textbf{Fluid Intelligence:} Ability to solve problems by coming up with new innovative ways to attack the problem. Peaks somewhere between ages 20-25 and then slowly declines afterward.

\textbf{Crystalized Intelligence:} Ability to solve problems from accumulated knowledge (make connections). This grows stronger with age.

The problem comes when people neglect making the jump from fluid intelligence to leverage crystallized intelligence during midlife. Important to note, these intelligences are not related to IQ.

From the book: "From Strength to Strength" by Arthur C. Brooks. (amazon link)
\end{minipage}
\end{center}
\vspace*{\fill}